\section{Carry a raw state across read boundaries}
Let \(U(r,B)\) mean updating raw state \(r\) with byte sequence \(B\). The update
contract requires
\[
  U(U(r,A),B)=U(r,A\Vert B),
\]
where \(A\Vert B\) is concatenation. The left side uses two calls but represents
one message. It neither reinitializes nor complements the state between the calls.
This is what lets a buffer of seven bytes process a file of any number of bytes.

\begin{table}[htbp]
  \centering
  \caption{Prefix states for the nine-byte check message.}
  \tagpdfsetup{table/header-rows={1}}
  \begin{tabular}{rll}
    \toprule
    Bytes consumed & Raw state         & Checksum if the message ended here \\
    \midrule
    0              & \texttt{ffffffff} & \texttt{00000000}                  \\
    1              & \texttt{7c231048} & \texttt{83dcefb7}                  \\
    3              & \texttt{77b79c2d} & \texttt{884863d2}                  \\
    6              & \texttt{f68d2c9e} & \texttt{0972d361}                  \\
    7              & \texttt{affc9660} & \texttt{5003699f}                  \\
    9              & \texttt{340bc6d9} & \texttt{cbf43926}                  \\
    \bottomrule
  \end{tabular}
\end{table}

With three-byte reads, the file reader retains the raw states in rows three, six and
nine. With a seven-byte first read, it retains the states in rows seven and nine.
Both routes end at the same state because they consumed the same ordered bytes.
The complemented column is useful for checking a prefix as an independent message;
those values are not the arguments for continuing the raw-state interface.

The two methods must also compose with each other. Process a prefix with the bit
routine and its suffix with the table routine, then reverse the methods. Carry the
actual returned state between calls. Checking only two complete-message results would
not establish that their intermediate state conventions agree.

\subsection*{Exercise: Explain a boundary-dependent failure}
A file checker succeeds when the entire file fits in one read but fails with a small
buffer. Trace the effect of complementing after each read, and then the effect of
reinitializing before each read. Explain which prefix information each mistake changes
or discards.
